\section{Results and Discussion}
\label{sec:results}

\subsection{Task 1: Text Generation}

\subsubsection*{Training Performance}

Table~\ref{tab:task1_training} reports wall-clock training time, peak memory usage,
final cross-entropy loss, and perplexity for all six text-generation models.
Perplexity is computed as $e^{\text{loss}}$ and indicates how many tokens the model
is effectively choosing between at each step — lower is better.

\begin{table}[H]
\centering
\caption{Text generation model training performance (20 epochs, CPU).}
\label{tab:task1_training}
\begin{tabular}{llrrrr}
\toprule
\textbf{Model} & \textbf{Level} & \textbf{Time (s)} & \textbf{Memory (MB)}$^*$ & \textbf{Loss} & \textbf{Perplexity} \\
\midrule
Markov Chain & Char & 0.15  & 2.5   & —      & —     \\
Markov Chain & Word & 0.01  & 2.0   & —      & —     \\
RNN          & Char & 101.6 & 828.1 & 0.6689 & 1.95  \\
RNN          & Word & 18.4  & 2.0   & 0.7746 & 2.17  \\
LSTM         & Char & 107.6 & 2.5   & \textbf{0.2943} & \textbf{1.34} \\
LSTM         & Word & 20.2  & 0.0   & 0.7861 & 2.19  \\
\bottomrule
\end{tabular}
\smallskip

\raggedright\footnotesize
$^*$ Memory is measured as the process RSS delta via \texttt{psutil} and is not
a reliable indicator of model size. The 828.1 MB spike for Char RNN is almost
certainly a measurement artefact: RSS captures OS-level page allocations, which
can include gradient buffers, dataloader workers, and OS-level fragmentation that
persist across models in the same process. All four neural models have
fewer than 500K parameters and should consume well under 10 MB of actual weights.
These figures should not be used to compare memory efficiency.
\end{table}

\subsubsection*{Discussion}

\textbf{Char LSTM achieves the lowest loss (0.2943, perplexity 1.34)} of all
text-generation models. A perplexity of 1.34 means the model is on average
uncertain between only 1.34 tokens at each step. Character sequences have high
local regularity (spelling patterns, common suffixes) that LSTM gates can exploit
by retaining memory over long windows. RNN achieves perplexity 1.95 at the same
level, directly illustrating the practical cost of the vanishing gradient problem.

\textbf{Word-level models converge more slowly} than their character-level
counterparts. On a corpus of $\approx$12,000 words, each word-level token is seen
far fewer times than each character, giving the model less signal per epoch. The
Word LSTM (perplexity 2.19) performs slightly worse than the Word RNN (2.17) —
the LSTM's larger parameter count requires more data to fit well within 20 epochs.

\textbf{Markov Chain models} train in under 1 second but offer no meaningful loss
or perplexity because they are not optimised by gradient descent.

\subsubsection*{Qualitative Output Comparison — Character-Level}

Table~\ref{tab:task1_char_outputs} shows character-level outputs for two prompts.
The progression from Markov Chain to LSTM illustrates increasing coherence.

\begin{table}[H]
\centering
\caption{Sample character-level outputs (truncated to 80 chars).}
\label{tab:task1_char_outputs}
\renewcommand{\arraystretch}{1.3}
\begin{tabular}{p{3cm} p{3.5cm} p{3.5cm} p{3.5cm}}
\toprule
\textbf{Prompt} & \textbf{Markov Chain} & \textbf{RNN} & \textbf{LSTM} \\
\midrule
``The ancient forest'' &
The analyzes the melody impressive principles\ldots &
The ancient forestery and the dish in many ways\ldots &
The ancient foresthes generously among the elegant perspectives\ldots \\
\midrule
``Science reveals'' &
ScienPeople have always been associated without\ldots &
Science reveals continue to evolve. Meanwhile\ldots &
Science reveals the boundaries of reliable taches\ldots \\
\bottomrule
\end{tabular}
\end{table}

\subsubsection*{Qualitative Output Comparison — Word-Level}

Table~\ref{tab:task1_word_outputs} shows word-level outputs for the same prompts.
All three models produce real words, but the Markov Chain tends to copy memorised
phrases, while the neural models generalise beyond the training n-grams — though
they introduce \texttt{<UNK>} tokens when the prompt contains words not seen
during training.

\begin{table}[H]
\centering
\caption{Sample word-level outputs (truncated to 80 chars).}
\label{tab:task1_word_outputs}
\renewcommand{\arraystretch}{1.3}
\begin{tabular}{p{3.2cm} p{3.3cm} p{3.3cm} p{3.3cm}}
\toprule
\textbf{Prompt} & \textbf{Markov Chain} & \textbf{RNN} & \textbf{LSTM} \\
\midrule
``the ancient forest stretches'' &
the ancient records play a vital role in progress\ldots &
the ancient forest stretches toward new strategies\ldots &
the ancient forest stretches freely. when the open meadow grows\ldots \\
\midrule
``the computer processes data'' &
the computer processes the boundaries of aromatic wines\ldots &
the computer processes \texttt{<UNK>} in recent years, the chef presents\ldots &
the computer processes \texttt{<UNK>} the methods that make life inclusive\ldots \\
\bottomrule
\end{tabular}
\end{table}

The \texttt{<UNK>} tokens in the word-level RNN and LSTM outputs highlight a key
limitation of word-level tokenization: any word in the prompt that falls below the
vocabulary frequency threshold is replaced with a generic unknown token, breaking
the prompt's continuity. The Markov Chain avoids this because it falls back to a
random known n-gram rather than substituting a token.

\subsection{Task 2: Chatbots}

\subsubsection*{Training Performance}

Table~\ref{tab:task2_training} compares the two chatbot models.

\begin{table}[H]
\centering
\caption{Chatbot model training performance (30 epochs, CPU).}
\label{tab:task2_training}
\begin{tabular}{lrrrr}
\toprule
\textbf{Model} & \textbf{Parameters} & \textbf{Time (s)} & \textbf{Loss} & \textbf{Perplexity} \\
\midrule
LSTM Chatbot        & 6,019,947 & 105.2 & \textbf{0.1014} & \textbf{1.11} \\
Transformer Chatbot & 1,971,819 & 33.4  & 4.3724          & 79.23         \\
\bottomrule
\end{tabular}
\end{table}

The perplexity gap is striking: the LSTM is uncertain between only 1.11 tokens per
step, while the Transformer is uncertain between 79 tokens — effectively sampling
from a large portion of the vocabulary at each decoding step.

\subsubsection*{Loss Curves}

Figure~\ref{fig:loss_curves} shows the training loss curves for both chatbot models
over 30 epochs.

\begin{figure}[H]
\centering
\includegraphics[width=0.85\linewidth]{figures/chatbot_loss_curves.png}
\caption{Training loss curves for the LSTM and Transformer chatbots (30 epochs).}
\label{fig:loss_curves}
\end{figure}

\subsubsection*{Response Quality}

Table~\ref{tab:task2_outputs} shows responses to representative test questions.

\begin{table}[H]
\centering
\caption{Sample chatbot responses to test questions.}
\label{tab:task2_outputs}
\renewcommand{\arraystretch}{1.3}
\begin{tabular}{p{4cm} p{5cm} p{5cm}}
\toprule
\textbf{Question} & \textbf{LSTM Response} & \textbf{Transformer Response} \\
\midrule
What is the capital of France? &
\textit{the capital of france is paris.} &
\textit{the capital of course. the capital of course. the capital of course.} \\
\midrule
How does gravity work? &
\textit{gravity works by creating an attractive force proportional to the masses\ldots} &
\textit{good question. the answer: the answer: the answer: the answer\ldots} \\
\midrule
What is photosynthesis? &
\textit{photosynthesis is the process by which plants convert sunlight into energy.} &
\textit{of course. the capital of course. the capital of course.} \\
\midrule
Hello! &
\textit{hello! how can i help you today?} &
\textit{to the process of course. the process of course.} \\
\bottomrule
\end{tabular}
\end{table}

\subsubsection*{Discussion}

\textbf{The LSTM chatbot achieves a final loss of 0.1014} and produces factually
correct, coherent answers on nearly all 10 test questions. Two factors explain this:
(i) it has 3$\times$ more parameters (6M vs.\ 2M), giving it greater capacity to
memorise the 2,000 QA pairs; and (ii) teacher forcing provides a strong per-step
training signal throughout training.

\textbf{The Transformer chatbot achieves a final loss of 4.3724 (perplexity 79.23)}
and exhibits \textit{repetition collapse} — repeatedly generating phrases such as
\textit{``the capital of course.''} regardless of the input. This is consistent with
an under-trained model whose output distribution has collapsed onto a small set of
high-frequency tokens.

The explanation offered here — insufficient data and epochs combined with a
conservative warmup LR schedule — is plausible but not confirmed by this experiment.
Confirming it would require at minimum:
\begin{itemize}
    \item \textbf{Extended training run} (e.g., 100--200 epochs) to observe whether
          the Transformer loss does in fact decrease past its current plateau.
    \item \textbf{Learning curve inspection} across all epochs to determine whether
          the model is still improving at epoch 30 or has genuinely stalled.
    \item \textbf{LR schedule ablation}: comparing the warmup schedule against a
          fixed or cosine-annealed rate to isolate the schedule's contribution.
\end{itemize}
These experiments were not conducted within the scope of this project, and the
convergence failure should therefore be treated as an observation rather than a
diagnosed cause.

\textbf{The Transformer trains 3$\times$ faster per epoch} (33s vs.\ 105s) due to
full sequence parallelism — self-attention processes all tokens simultaneously,
unlike the LSTM which must process them one at a time.

\subsection{Model Complexity vs.\ Performance}

Table~\ref{tab:complexity} brings all six neural models together to examine whether
parameter count predicts performance. Markov Chain models are excluded as they have
no learnable parameters.

\begin{table}[H]
\centering
\caption{Parameter count vs.\ final performance for all neural models.}
\label{tab:complexity}
\begin{tabular}{llrrr}
\toprule
\textbf{Model} & \textbf{Task} & \textbf{Parameters} & \textbf{Loss} & \textbf{Perplexity} \\
\midrule
Char RNN            & Text gen  &    70,980 & 0.6689 &  1.95 \\
Char LSTM           & Text gen  &   244,548 & \textbf{0.2943} & \textbf{1.34} \\
Word RNN            & Text gen  &   289,456 & 0.7746 &  2.17 \\
Word LSTM           & Text gen  &   463,024 & 0.7861 &  2.19 \\
Transformer Chatbot & Chatbot   & 1,971,819 & 4.3724 & 79.23 \\
LSTM Chatbot        & Chatbot   & 6,019,947 & 0.1014 &  1.11 \\
\bottomrule
\end{tabular}
\end{table}

Two observations stand out:

\begin{itemize}
    \item \textbf{More parameters do not guarantee better performance.} Char LSTM
    (245K) achieves a lower loss than Word LSTM (463K), and the Transformer (2M)
    performs far worse than the LSTM Chatbot (6M). In both cases the gap is
    explained by training dynamics — data volume, training duration, and LR
    schedule — not architecture alone.

    \item \textbf{Task difficulty dominates scale.} The chatbot task
    (sequence-to-sequence over QA pairs) is substantially harder than next-token
    prediction over a fixed corpus. Loss values across the two tasks are not
    directly comparable; perplexity within each task is the meaningful metric.
\end{itemize}
